\documentclass[12pt]{article}
\usepackage[utf8]{inputenc}
\usepackage[T1]{fontenc}
\usepackage{hyperref}
\title{Engenni Language Resources for K-12 Curriculum Development}
\author{Compiled by Claude Sonnet 5.5 (claude-sonnet-5-5) for Andr\'as Kornai}
\date{}
\begin{document}
\maketitle
\section{Introduction}
Engenni (Egene; ISO 639-3 enn, Glottocode enge1239) is a Niger-Congo language of Nigeria. Speaker estimates in the sources differ widely (about 20,000 to over 200,000). Very little educational material was found.
\section{Primary Educational Resources}
No primers or textbooks were found in the search results.
\section{Linguistic and Reference Materials}
\begin{itemize}
\item Wikipedia overview, which mentions an orthography revised in 2011 (10 vowels): \url{https://en.wikipedia.org/wiki/Engenni_language}
\item Omniglot page on the Egene writing system: \url{https://www.omniglot.com/writing/egene.htm}
\item Glottolog reference entry for Elaine Thomas (1974), a description of Engenni: \url{https://glottolog.org/resource/reference/id/85655}
\end{itemize}
\section{Religious and Community Resources}
Joshua Project lists Bible portions and a 1977 New Testament, with no complete Bible reported: \url{https://joshuaproject.net/languages/enn}
\section{Digital and Audio Resources}
None were found in this search.
\section{Curriculum Development Context}
The language is not known to be taught in schools. Any curriculum would first need the orthography and its materials checked with the community; no dictionary was found.
\end{document}